\section{Using comments in \LaTeX}
	In \LaTeX, \emph{comments} allow authors to insert non-printing notes, suppress undesired whitespace, or temporarily disable sections of code during document development. Comments are ignored by the \LaTeX\ compiler and do not appear in the generated output document.
	
	A comment is introduced using the percent character (\%). Upon encountering a \% symbol in standard text mode, \LaTeX\ ignores the \% itself and everything following it on that same line until the end-of-line character.
	
	For instance, the following \LaTeX\ code snippet:
\lstinputlisting{codes/snippet-comments.tex}
	produces the following output:
	\lstoutput{codes/snippet-comments.tex}
	
	\subsection{Utility of comments}
		As in any programming language, comments in \LaTeX\ have numerous utilities, some of which are outlined below.
		\begin{enumerate}
			\item Authors can leave descriptive annotations, structural markers, or to-do notes for themselves or co-authors without affecting the typeset document. \LaTeX\ templates frequently use comments to guide users on required preamble setups or placeholder entries.
			
			\item Rather than deleting commands or text paragraphs when testing alternative layouts, authors can ``comment out'' lines by prefixing them with \%. This preserves the original code, allowing it to be easily restored by simply removing the percent sign.
			
			\item Because a single line break in source code acts as a space character, placing a \% at the very end of a line (immediately following the last character without an intervening space) masks the newline. This technique is widely used when constructing macros, defining custom environments, or placing inline elements side-by-side where extra whitespace would disrupt alignment or page layout.
		\end{enumerate}
	
	